% 1 Slide for program flow
\subsection{Program Flow}\label{subsection:program_flow}
\begin{frame}{Program Flow}
    \begin{tikzpicture}
        \node[draw] (main) {main.py};
        \node[draw, below=of main] (calibrator) {ball\_detector\_calibration.py};
        \node[draw, above=of main] (camera) {camera\_handler.py};
        \node[draw, right=of main] (detector) {ball\_detector.py};
        \node[draw, right=of detector] (output) {output.py};
        \draw[->] (main.south) -- (calibrator.north) node[midway, right] {Called before first call of detector};
        \draw[<->] (main.north) -- (camera.south) node[midway, right] {Abstraction of picture source};
        \draw[->] (main.east) -- (detector.west);
        \draw[->] (detector.east) -- (output.west);
    \end{tikzpicture}
\end{frame}

% 1-2 Slides for corner detection
\subsection{Corner Detection}\label{subsection:program_cornerDetection}
\begin{frame}{Corner Detection}
    \begin{itemize}
        \item Algorithm for top-left corner, but applies to all corners when adjusted for direction
        \begin{itemize}
            \item Morphological opening used to clean greyscale picture before algorithm
            \item Use Manhattan-distance to find the closest pixel
        \end{itemize}
        \includegraphics[width=0.25\textwidth]{chapters/path\_of\_corner\_detection.png}
        \begin{enumerate}
            \item Find top-left corner of support beam in greyscale picture
            \item Find top-left corner of wall in greyscale picture
            \item Find top-left corner of field in selection of green colour
        \end{enumerate}
        \item Perform manual correction, in case automatic detection fails
        \item Use \texttt{\href{https://docs.opencv.org/3.4.20/da/d54/group__imgproc__transform.html\#ga8c1ae0e3589a9d77fffc962c49b22043}{cv2.getPerspectiveTransform()}} to calculate projection matrix from field corners
    \end{itemize}
\end{frame}

\subsection{Challenges with corner detection}\label{subsection:program_challenges}
\begin{frame}{Challenges with corner detection}
    \begin{itemize}
        \item During game
        \begin{itemize}
            \item Hands above or in the field.
            \item Low FPS
        \end{itemize}
        \item During calibration
        \begin{itemize}
            \item Different lighting conditions are very effective at throwing off greyscale thresholding.
            \item Bright objects (and people) surrounding the table can be detected as corners.
            \item The silver rods throw off brightness based corner detection, prevent morphological closing and opening.
        \end{itemize}
        \item In general
        \begin{itemize}
            \item The camera was moved during development. Test pictures taken for development before the camera was moved are no longer accurate.
        \end{itemize}
    \end{itemize}
\end{frame}

% 1 Slide for slide perspective transformation + coordinate calculation
\subsection{Perspective Transformation}\label{subsection:program_perspctiveTransformation}
\begin{frame}{Perspective Transformation}
    \begin{enumerate}
        \item Select orange area from HSV picture using \texttt{\href{https://docs.opencv.org/3.4.20/d2/de8/group__core__array.html\#ga48af0ab51e36436c5d04340e036ce981}{cv2.inRange()}}
        \item Check size of area
        \begin{itemize}
            \item size is zero: No ball on field
            \item normal: Ball found
            \item very large: Hand over field
            \item \texttt{\href{https://docs.opencv.org/3.4.20/d3/dc0/group__imgproc__shape.html\#ga556a180f43cab22649c23ada36a8a139}{M = cv2.moments}}
            \item \texttt{M["m00"]} provides size
        \end{itemize}
        \item Calculate center of orange area
        \begin{itemize}
            \item \texttt{int(M["m10"] / M["m00"])} to get x-coordinate
            \item \texttt{int(M["m01"] / M["m00"])} to get y-coordinate
            \item Future Idea: Detect multiple balls
        \end{itemize}
        \item Project center cordinates from image space to field space
        \begin{itemize}
            \item Use previously calculated perspective projection matrix
            \item Apply Translation from top-left alignment to center alignment
        \end{itemize}
    \end{enumerate}
\end{frame}